\section{Introduction}
\label{sec:intro}

A system that could recover the meaning of silently imagined speech from scalp EEG would be useful
to people who cannot speak, and it would be cheap and non-invasive. Reported accuracies for this
task are low, so small differences carry the argument: a result of 14\% on a 39-way problem is
interesting only if one is confident that the 11 points above chance come from the imagined
sentence and not from something else the recording happens to encode.

EEG makes that confidence hard to earn. Recordings are long, and EEG contains temporal
correlations that make trials recorded close together resemble each other regardless of the
stimulus; a classifier can exploit them by learning \emph{when} a trial was recorded
\citep{li2021perils}. Whether that matters depends on
two design choices that are easy to make and easy to overlook: the order in which stimuli are
presented, and how test trials are drawn from the recording.

We examine one benchmark where both choices interact with a third, the rule used to report
accuracy. Chisco \citep{zhang2024chisco} is a large public imagined-speech corpus: over 6{,}000
sentences per participant, recorded across five days and labelled with 39 semantic categories.
Its authors report 13.88\% category accuracy for participant S1 against a 5.52\% shuffled-label
baseline, and conclude that the EEG contains semantic information. The released code assigns every
fifth trial in recording order to the evaluation set and reports the best evaluation accuracy seen
during training.

Our central observation needs no EEG. In the released recording order, consecutive trials share a
category 19.7\% of the time, against 2.95\% if the same categories were shuffled. On the released
evaluation set, a classifier that copies the label of the preceding training trial scores
\textbf{18.8\%}, above the published EEG result. The EEG, in turn, carries the time stamp this
classifier needs: per-trial channel means and variances identify the recording run of a held-out
trial with 98.8\% accuracy. Running the released script with its documented command on our
hardware prints 9.69\%, not 13.88\%. Under the same rule, a model trained on permuted labels still
reports 6.4\%, and real labels add 2.4 points to that, not the 8.4 the paper reports. When whole
runs are held out so that evaluation trials have no neighbours in the training set, the margin
between real and permuted labels falls to 0.64 points, within the variation between seeds.

Our contributions are:
\begin{itemize}
\item An audit of the released 39-way benchmark on its exact evaluation split: trial position alone
predicts the category better than the published EEG model, and content-free EEG summary statistics
carry enough recording-time information to exploit that structure (Section~\ref{sec:audit}).
\item A re-run of the published configuration, including the released script itself, that reads
three selection rules from the same training runs and compares each with a permuted-label null
(Section~\ref{sec:results}). The published best-checkpoint rule adds 1.4--1.6 points over a
validation-selected checkpoint and reports 5.2--6.4\% for models trained without label information.
\item A split comparison showing the real-vs-permuted margin shrinking from 2.11 points on the
released split to 0.64 under run holdout (Section~\ref{sec:splits}).
\item A short protocol for temporally ordered decoding benchmarks: report a position-only baseline
on the exact evaluation split, hold out whole runs or sessions, select checkpoints on data that is
never scored, and compare against a permuted-label null run through the same pipeline.
\end{itemize}

We do not claim that the Chisco EEG contains no semantic information, or that imagined-speech
decoding from EEG is infeasible. Our claims are about what this benchmark's protocol can and
cannot show, for one participant.
